\documentclass[10pt, a4paper]{article}
\usepackage{preamble}

\title{Probability II \\
    \large Prereading Notes}
\author{Luke Phillips}
\date{January 2026}

\begin{document}

\maketitle

\newpage

\tableofcontents

\newpage

\section{Probability and independence}

\subsection{Probability spaces}

\begin{definition}[Sigma-field]
    Let $\mathcal{F}$ be a collection of subsets of $\Omega$.
    We call $\mathcal{F}$ a $\sigma$-field or a $\sigma$-algebra if it has the following properties:
    \begin{enumerate}[label = \arabic*.]
        \item $\emptyset \in \mathcal{F}$.

        \item If $A \in \mathcal{F}$,
        then $A ^ c \in \mathcal{F}$.

        \item If $A_1, A_2, \dotsc \in \mathcal{F}$,
        then $\bigcup_{k = 1}^{\infty}A_k \in \mathcal{F}$.
    \end{enumerate}
\end{definition}

\begin{definition}[Probability]
    Let $\Omega$ be a sample space,
    and $\mathcal{F}$ be a $\sigma$-field of events in $\Omega$.
    A probability distribution $\P$ on $(\Omega, \mathcal{F})$ is a collection of numbers $\P(A)$,
    $A \in \mathcal{F}$,
    possessing the following properties:
    \begin{enumerate}[label = A\arabic*]
        \item For every event $A \in \mathcal{F}$,
        $\P(A) \geq 0$.

        \item $\P(\Omega) = 1$.

        \item For any pair of incompatible events $A$ and $B$
        (i.e.,
        $A \cap B = \emptyset$),
        $\P(A \cup B) = \P(A) + \P(B)$.

        \item For any countable collection $A_1, A_2, \dotsc$ of incompatible events
        (mutually disjoint)
        \[
        \P\left(\bigcup_{k = 1}^{\infty}A_k\right) = \sum_{k = 1}^{\infty}\P(A_k).
        \]
    \end{enumerate}
\end{definition}

\begin{definition}
    A probability space is a triple $(\Omega, \mathcal{F}, \P)$,
    where $\Omega$ is a sample space,
    $\mathcal{F}$ is a $\sigma$-field of events in $\Omega$,
    and $\P(\cdot)$ is a probability measure on $(\Omega, \mathcal{F})$.
\end{definition}

\subsubsection{Generated \texorpdfstring{$\sigma$}{}-fields}
Start by checking that an intersection of two $\sigma$-fields is a $\sigma$-field.
So check that it satisfies all of the properties to be a sigma field.

\begin{definition}
    Let $\mathcal{X}$ be an arbitrary set and let $\mathcal{D}$ be an arbitrary collection of subsets of $\mathcal{X}$.
    Let $\mathcal{F}_{\alpha}$,
    $\alpha \in \mathcal{A}$,
    be the collection of all sigma-fields in $\mathcal{X}$ which contain $\mathcal{D}$,
    namely $\mathcal{D} \subset \mathcal{F}_{\alpha}$ for all $\alpha \in \mathcal{A}$.
    Then $\mathcal{G} \coloneqq \mathcal{G}_{\mathcal{D}} = \bigcap_{\alpha \in \mathcal{A}}\mathcal{F}_{\alpha}$ is the smallest sigma-field of subsets of $\mathcal{X}$ containing $\mathcal{D}$,
    also known as the $\sigma$-field generated by $\mathcal{D}$.
\end{definition}

\subsection{Monotone sequences of events}

\begin{definition}
    A sequence $(A_n)_{n \geq 1}$ of events is increasing if $A_n \subseteq A_{n + 1}$ for all $n \geq 1$.
    It is decreasing if $A_n \supseteq A_{n + 1}$ for all $n \geq 1$.
\end{definition}

\begin{lemma}
    If $(A_n)_{n \geq 1}$ is increasing with $A \coloneqq \lim_{n}A_n \equiv \bigcup_{n \geq 1}A_n$,
    then
    \[
    \P(A) = \P\left(\lim_{n \to \infty}A_n\right) = \lim_{n \to \infty}\P(A_n).
    \]
    If $(A_n)_{n \geq 1}$ is a decreasing sequence with $A \coloneqq \lim_{n}A_n \equiv \bigcap_{n \geq 1}A_n$,
    then
    \[
    \P(A) = \P\left(\lim_{n \to \infty}A_n\right) = \lim_{n \to \infty}\P(A_n).
    \]

    \begin{proof}
        Let $(A_n)_{n \geq 1}$ be increasing with $A = \bigcup_{n \geq 1}A_n$.
        Denote $C_1 = A_1$ and,
        for $n \geq 2$,
        put $C_n = A_n \setminus A_{n - 1} = A_n \cap A_{n - 1}^{c}$.
        We then have
        \[
        A_n = \bigcup_{k = 1}^{n}A_k = \bigcup_{k = 1}^{n}C_k,\qquad\bigcup_{k = 1}^{\infty}A_k = \bigcup_{k = 1}^{\infty}C_k.
        \]
        Since the events in $(C_k)_{k \geq 1}$ are mutually incompatible,
        the $\sigma$-additivity property of probability gives
        \[
        \P(A) = \P\left(\bigcup_{k \geq 1}A_k\right) = \P\left(\bigcup_{k \geq 1}C_k\right) = \sum_{k \geq 1}\P(C_k) \leq 1.
        \]
        Therefore
        \[
        0 \leq \P(A) - \P(A_n) = \P(A \setminus A_n) = \P\left(\bigcup_{k > n}C_k\right) = \sum_{k > n}\P(C_k) \to 0
        \]
        as $n \to \infty$,
        as a tail sum of a convergent series $\sum_{k \geq 1}\P(C_k)$.
        Similar argument for decreasing sequences.
    \end{proof}
\end{lemma}

\subsection{Independence of events}

\begin{definition}
    Let $(\Omega, \mathcal{F}, \P)$ be a probability space.
    Two events $A, B, \in \mathcal{F}$ are independent,
    if their joint probability factorises,
    \[
    \P(A \cap B) = \P(A)\P(B).
    \]
\end{definition}

\begin{definition}
    Let $(\Omega, \mathcal{F}, \P)$ be a probability space.
    A finite or infinite collection of events $(A_{\alpha})_{\alpha \in \mathcal{A}}$ is
    (mutually)
    independent,
    if the probability of any finite sub-collection factorises,
    namely,
    for all integer $k \geq 1$ and all $\alpha_1, \dotsc, \alpha_k \in \mathcal{A}$,
    \[
    \P\left(\bigcap_{\ell = 1}^{k}A_{\alpha_{\ell}}\right) = \prod_{\ell = 1}^{k}\P(A_{\alpha_{\ell}}).
    \]
\end{definition}

\subsection{Random variables}

\begin{lemma}
    Let $X$ be a finite random variable on a probability space $(\Omega, \mathcal{F}, \P)$,
    i.e.,
    $\P(|X| < \infty) = 1$.
    Its cumulative distribution function $F_{X}(y) \coloneqq \P(X \leq y)$ has the following properties:
    \begin{enumerate}[label = (\alph*)]
        \item $F_{X}(y)$ is a non-decreasing function of $y$ such that $\lim_{y \to -\infty ^ {+}}F_{X}(y) = 0$ and $\lim_{y \to \infty ^ {-}}F_{X}(y) = 1$.

        \item $F_{X}(y)$ is a right-continuous function of $y$,
        i.e.,
        $\lim_{\varepsilon \to 0 ^ {+}}F_{X}(y + \varepsilon) = F_{X}(y)$ for all $y \in \R$.
    \end{enumerate}
\end{lemma}

\subsubsection{Independent random variables}
\begin{definition}
    Two random variables $X, Y$ on the same probability space $(\Omega, \mathcal{F}, \P)$ are independent,
    if
    \[
    \forall x, y \in \R,\qquad \P(X \leq x, Y \leq y) = \P(X \leq x)\P(Y \leq y).
    \]
\end{definition}

\begin{definition}
    An arbitrary collection of random variables $(X_{\alpha})_{\alpha \in \mathcal{A}}$ on some probability space $(\Omega, \mathcal{F}, \P)$ is
    (mutually)
    independent,
    if any finite sub-collection of these variables is independent,
    namely,
    for all integer $k \geq 1$,
    all $\alpha_1, \dotsc, \alpha_k \in \mathcal{A}$,
    and all real $x_1, \dotsc, x_k$,
    \[
    \P(X_{\alpha_1} \leq x_{\alpha_1}, \dotsc, X_{\alpha_k} \leq x_{\alpha_k}) = \prod_{\ell = 1}^{k}\P(X_{\alpha_{\ell}} \leq x_{\alpha_{\ell}}).
    \]
\end{definition}

\newpage

\section{General sequences of events}

\subsection{Two special limit events}
It is useful to describe convergence of events in terms of their indicators
\[
\mathbbm{1}_{A}(\omega) = \begin{cases}
    1, &\text{if }\omega \in A, \\
    0, &\text{if }\omega \notin A.
\end{cases}
\]

\subsection{Borel-Cantelli lemma}

We can define the event
"infinitely many of the events $A_n$ occur"
as
\[
\{A_n\text{ infinitely often}\} = \bigcap_{n \geq 1}\bigcup_{k = n}^{\infty}A_k.
\]

\begin{lemma}[Borel-Cantelli lemma]
    Let $A = \bigcap_{n \geq 1}\bigcup_{k = n}^{\infty}A_k$ be the event $\{A_n\text{ infinitely often}\}$.
    Then:
    \begin{enumerate}[label = (\alph*)]
        \item If $\sum_{k}\P(A_k) < \infty$,
        then $\P(A) = 0$,
        i.e.,
        with probability one only finitely many of the $A_k$ occur.

        \item If $\sum_{k}\P(A_k) = \infty$ and $A_1, A_2, \dotsc$ are independent events,
        then $\P(A) = 1$.
    \end{enumerate}

    \begin{proof}\phantom{}
        \begin{enumerate}[label = (\alph*)]
            \item For every $n \geq 1$,
            let $B_n \coloneqq \bigcup_{k = n}^{\infty}A_k$ be the event that at least one of $A_k$ with $k \geq n$ occurs.
            As $A \subset B_n$ for all $n \geq 1$,
            we have
            \[
            0 \leq \P(A) \leq \P(B_n) \leq \sum_{k = n}^{\infty}\P(A_k) \to 0
            \]
            as $n \to \infty$,
            whenever $\sum_{k}\P(A_k) < \infty$.
            Consequently,
            $\P(A) = 0$.

            \item The event $A ^ c = \{A_n\text{ occur finitely often}\}$ is related to the sequence
            \[
            B_n ^ c = \bigcap_{k = n}^{\infty}A_{k}^{c} \equiv \{\text{none of $A_k$,
            $k \geq n$,
            occurs}\}\quad\text{via}\quad A ^ c = \bigcup_{n}\bigcap_{k = n}^{\infty}A_{k}^{c} = \bigcup_{n}B_{n}^{c},
            \]
            so for $\P(A ^ c) = 0$ it is sufficient to show that $\P(B_{n}^{c}) = 0$ for all $n \geq 1$.
            By independence and the elementary inequality $1 - x \leq e ^ {-x}$ with $x \geq 0$,
            we get
            \[
            \P\left(\bigcap_{k = n}^{m}A_{k}^{c}\right) = \prod_{k = n}^{m}\P(A_{k}^{c}) = \prod_{k = n}^{m}(1 - \P(A_k)) \leq \exp\left\{-\sum_{k = n}^{m}\P(A_k)\right\}
            \]
            so that
            \[
            \P(B_{n}^{c}) = \lim_{m \to \infty}\P\left(\bigcap_{k = n}^{m}A_{k}^{c}\right) \leq \exp\left\{-\sum_{k = n}^{\infty}\P(A_k)\right\} = 0,
            \]
            as the sum diverges.
            Hence,
            $0 \leq \P(A ^ c) \leq \sum_{n \geq 1}\P(B_{n}^{c}) = 0$ by sub-additivity,
            implying that $\P(A) = 1$.
        \end{enumerate}
    \end{proof}
\end{lemma}

\newpage

\section{Convergence of random variables}

\subsection{Main modes of convergence}

\begin{definition}
    For a sequence of random variables $(X_n)_{n \geq 1}$ converges pointwise or surely to $X$ if $X_n(\omega) \to X(\omega)$ for all $\omega \in \Omega$ as $n \to \infty$.
\end{definition}

\subsubsection{Convergence in \texorpdfstring{$L ^ r$}{}}
\begin{definition}
    For fixed $r > 0$,
    let $X$ and $(X_n)_{n \geq 1}$ be random variables such that $\E(|X| ^ r)$ and all $\E(|X_n| ^ r)$ are finite.
    Then the sequence $X_n$ converges to $X$ in $L ^ r$ as $n \to \infty$,
    if $\E(|X_n - X| ^ r) \to 0$ as  as $n \to \infty$.
\end{definition}

\begin{example}
    Let $(X_n)_{n \geq 1}$ be a sequence of random variables such that for some real numbers $(a_n)_{n \geq 1}$,
    we have
    \[
    \P(X_n = a_n) = p_n,\qquad \P(X_n = 0) = 1 - p_n.
    \]

    $X_n \to 0$ in $L ^ r$ if and only if $\E(|X_n| ^ r) = |a_n| ^ rp_n$ as $n \to \infty$.
\end{example}

\subsubsection{Convergence in probability}
\begin{definition}
    We say a sequence $(X_n)_{n \geq 1}$ of random variables converges to a random variable $X$ in probability
    (write $X_n \overset{p}{\to} X$)
    as $n \to \infty$,
    if for every fixed $\varepsilon > 0$
    \[
    \P(|X_n - X| \geq \varepsilon) \to 0\text{ as } n \to \infty.
    \]
\end{definition}

\begin{example}
    Let random variables $(X_n)_{n \geq 1}$ be such that for some real $(a_n)_{n \geq 1}$,
    we have
    \[
    \P(X_n = a_n) = p_n,\qquad \P(X_n = 0) = 1 - p_n
    \]
    where $a_n \to 0$ as $n \to \infty$.
    Does $X_n \to 0$ as $n \to \infty$ in $L ^ r$ and does $X_n \to 0$ as $n \to \infty$ in probability?

    \begin{solution}
        $X_n \to 0$ as $n \to \infty$ in $L ^ r$ if and only if
        \[
        \E(|X_n| ^ r) = |a_n| ^ rp_n \to 0\text{ as } n \to \infty
        \]
        but if $a_n \to 0$ then $|a_n| \to 0$ and hence $|a_n| ^ r \to 0$ so we get that it converges in $L ^ r$ for $r > 0$.
    \end{solution}
\end{example}

\begin{lemma}[General Markov Inequality]
    Let $g : [0, \infty) \to [0, \infty)$ be an increasing function.
    Then for each random variable $X \geq 0$ and a fixed constant $a > 0$ we have
    \[
    \P(X \geq a) \leq \frac{\E(g(X))}{g(a)}.
    \]

    \begin{proof}
        Since $g$ increases,
        for each fixed $a > 0$ and all $x \geq 0$ we have $g(x) \geq g(x)\mathbbm{1}_{x \geq a} \geq g(a)\mathbbm{1}_{x \geq a} \geq 0$.
        By applying this inequality to the random variable $X \geq 0$ and taking expectations,
        we obtain
        \[
        \E(g(X)) \geq g(a)\E(\mathbbm{1}_{X \geq a}) \equiv g(a)\P(X \geq a).
        \]
        As $g(a) > 0$,
        this is equivalent to the required inequality.
    \end{proof}
\end{lemma}

\begin{lemma}
    Let $(X_n)_{n \geq 1}$ be a sequence of random variables.
    If $X_n \overset{L ^ r}{\to} X$ for some fixed $r > 0$,
    then $X_n \overset{p}{\to} X$ as $n \to \infty$.

    \begin{proof}
        Applying the general Markov inequality to $g(x) = x ^ r$ and the random variable $|X_n - X| \geq 0$.
        Then,
        for every fixed $\varepsilon > 0$,
        \[
        \P(|X_n - X| \geq \varepsilon) \leq \frac{\E(|X_n - X| ^ r)}{\varepsilon ^ r} \to 0\text{ as } n \to \infty.
        \]
    \end{proof}
\end{lemma}

\subsubsection{Almost sure convergence}
\begin{definition}
    A sequence $(X_k)_{k \geq 1}$ of random variables in $(\Omega, \mathcal{F}, \P)$ converges,
    as $n \to \infty$,
    to a random variables $X$ with probability one
    (or almost surely)
    if
    \[
    \P\left(\{\omega \in \Omega\,:\, X_n(\omega) \to X(\omega)\text{ as } n \to \infty\}\right) = 1.
    \]
\end{definition}

\begin{lemma}
    Let $X$ and $(X_n)_{n \geq 1}$ be random variables.
    If,
    for every $\varepsilon > 0$,
    \[
    \sum_{n = 1}^{\infty}\P(|X_n - X| > \varepsilon) < \infty,
    \]
    then $X_n$ converges to $X$ almost surely.

    \begin{proof}
        Fix $\varepsilon > 0$ and let $A_n(\varepsilon) = \{\omega \in \Omega\,:\, |X_n(\omega) - X(\omega)| > \varepsilon\}$.
        By assumption,
        $\sum_{n}\P(A_n(\varepsilon)) < \infty$,
        and the Borel-Cantelli lemma,
        only finitely many $A_k(\varepsilon)$ occur with probability $1$.
        This means that for every fixed $\varepsilon > 0$ the event
        \[
        A(\varepsilon) \coloneqq \{\omega \in \Omega\,:\,|X_n(\omega) - X(\omega)| \leq \varepsilon\text{ for all $n$ large enough}\}
        \]
        has probability one.
        By monotonicity
        ($A(\varepsilon_1) \subset A(\varepsilon_2)$ if $\varepsilon_1 < \varepsilon_2$),
        the event
        \[
        \{\omega \in \Omega\,:\,X_n(\omega) \to X(\omega)\text{ as } n \to \infty\} = \bigcap_{\varepsilon > 0}A(\varepsilon) = \bigcap_{m \geq 1}A(1 / m)
        \]
        has probability one.
        The claim follows.
    \end{proof}
\end{lemma}

\end{document}